% Generated draft: requires mathematical and source review.
\documentclass[11pt]{article}
\usepackage[margin=25mm]{geometry}
\usepackage{fontspec}
\setmainfont{Latin Modern Roman}
\usepackage{amsmath,amssymb,mathtools,unicode-math}
\setmathfont{Latin Modern Math}
\usepackage{enumitem,needspace,xurl,hyperref}
\setlength{\parindent}{0pt}
\setlength{\parskip}{5pt}
\setlength{\emergencystretch}{3em}
\newcommand{\N}{\mathbb{N}}
\newcommand{\Z}{\mathbb{Z}}
\newcommand{\Q}{\mathbb{Q}}
\newcommand{\R}{\mathbb{R}}
\newcommand{\C}{\mathbb{C}}
\newcommand{\F}{\mathbb{F}}
\newcommand{\A}{\mathbb{A}}
\newcommand{\PP}{\mathbb{P}}
\newcommand{\E}{\mathbb{E}}
\newcommand{\eps}{\varepsilon}
\newcommand{\dd}{\,\mathrm{d}}
\begin{document}
\section*{UnsolvedMath ID: 30001498}
\textbf{Source problem code:} catalogue entry
\subsection{Definitions and mathematical statement}
For clarity, the following is the conventional cohomological definition; the supplied record does not specify the intended geometric hypotheses. Let \(k\) be a field, let \(X\) be a projective \(k\)-scheme, let \(L\) be a line bundle on \(X\), and let \(q\in\Z_{\geq 0}\). The line bundle \(L\) is \emph{\(q\)-ample} if, for every coherent sheaf \(\mathcal F\) on \(X\), there exists an integer \(m_0=m_0(\mathcal F,L,q)\geq 0\) such that
\[
H^i\!\left(X,\mathcal F\otimes L^{\otimes m}\right)=0
\qquad\text{for every integer }m\geq m_0\text{ and every integer }i>q.
\]
For \(q=0\), this is the usual cohomological characterization of ampleness. If \(q\geq\dim X\), the condition is automatic in this framework.

Numerical equivalence of line bundles means equality of their degrees on every integral curve in \(X\). This explains one conventional meaning of numerical data; the supplied record does not specify which intersection numbers, cycle cones, or other numerical data a characterization may use.

\textbf{Problem.} Give a necessary-and-sufficient numerical characterization of \(q\)-ampleness that specializes, when \(q=0\), to Kleiman's ampleness criterion in its applicable geometric setting.

\textbf{Unresolved scope.} The projective framework above supplies a definition, not an attribution of missing hypotheses to the source. The record does not specify the base field or its characteristic, whether \(X\) must be a variety, smooth, or irreducible, or the intended range of \(q\). It also gives no proposed numerical criterion. These choices require source clarification before a fully scoped problem statement can be asserted.
\subsection{Short English statement}
Find numerical conditions that are both necessary and sufficient for a line bundle to be \(q\)-ample, recovering Kleiman's criterion for ordinary ampleness at \(q=0\). The supplied record leaves the geometric hypotheses and permitted numerical data unspecified; the displayed definition uses a conventional projective framework.
\subsection{Sources}
\begin{itemize}
\item Local UnsolvedMath registry, supplied record with permanent ID \(30001498\), source problem code \texttt{catalogue entry}, title \emph{Numerical Characterization of \(q\)-Ampleness}, in the Oberwolfach Reports Open Problems collection.
\item \emph{Classical Algebraic Geometry} (2010), Oberwolfach Reports, report 27/2010, pp.~1573--1623, DOI \texttt{10.4171/owr/2010/27}: \url{https://doi.org/10.4171/owr/2010/27}. This is the report citation supplied by the registry; no precise problem page or named proposer is supplied.
\item Supplied literature source, labeled ``original OWR Report 27/2010'': \url{https://doi.org/10.4171/owr/2010/27}.
\item Literature assessment dated \texttt{2026-08-21}: the supplied curation labels the problem open, states that the report explicitly presents a numerical characterization of \(q\)-ampleness as a natural open problem, and describes available criteria as partial rather than a general Kleiman-type characterization. This is supplied dated curation, not an independently verified literature assessment.
\item The registry labels the statement \texttt{corrected\_verified} and records a manual semantic, typographic, and LaTeX review of the cleanup. These are supplied review claims and do not establish independent verification of the report, the mathematical scope, or the current problem status.
\end{itemize}
\end{document}
